\documentclass[12pt,oneside]{book}

% =====================================================
% METADATOS
% =====================================================
\newcommand{\universidad}{Universidad de Buenos Aires}
\newcommand{\facultad}{Facultad de Derecho}
\newcommand{\carrera}{}
\newcommand{\materia}{Derecho Procesal Civil y Comercial}
\newcommand{\titulo}{Unidad 21 -- Ejecuci\'on de Sentencia}
\newcommand{\autor}{Isaac Edgar Camacho Ocampo}
\newcommand{\anio}{2025}

% =====================================================
% PAQUETES
% =====================================================
\usepackage[utf8]{inputenc}
\usepackage[T1]{fontenc}
\usepackage[spanish,es-nodecimaldot]{babel}

\usepackage[
    top=2.5cm,
    bottom=2.5cm,
    left=3cm,
    right=2.5cm,
    headheight=15pt
]{geometry}

\usepackage{amsmath}
\usepackage{amssymb}
\usepackage{mathtools}

\usepackage{graphicx}
\usepackage{float}
\usepackage{caption}
\usepackage{subcaption}

\usepackage{booktabs}
\usepackage{array}
\usepackage{longtable}
\usepackage{tabularx}

\usepackage{enumitem}

\usepackage{xcolor}
\usepackage[most]{tcolorbox}

\usepackage{fancyhdr}
\usepackage{ragged2e}

\graphicspath{
    {./img/}
    {./graficos/}
    {./figuras/}
}

% =====================================================
% CONFIGURACI\'ON
% =====================================================
\setcounter{tocdepth}{2}

\setlength{\parindent}{0pt}
\setlength{\parskip}{0.5em}

\setlist[itemize]{
    topsep=0.3em,
    itemsep=0.2em
}

\setlist[enumerate]{
    topsep=0.3em,
    itemsep=0.2em
}

\clubpenalty=10000
\widowpenalty=10000

% =====================================================
% COMANDOS MATEM\'ATICOS
% =====================================================
\newcommand{\abs}[1]{\left|#1\right|}
\newcommand{\norm}[1]{\left\lVert#1\right\rVert}
\newcommand{\vect}[1]{\mathbf{#1}}
\newcommand{\deriv}[2]{\frac{d#1}{d#2}}
\newcommand{\pderiv}[2]{\frac{\partial#1}{\partial#2}}

% =====================================================
% ENTORNOS / CAJAS ACAD\'EMICAS
% =====================================================
\newtcolorbox{definition}[1]{
    enhanced,
    breakable,
    title={Definici\'on: #1},
    fonttitle=\bfseries,
    colback=blue!3,
    colframe=blue!50!black,
    boxrule=0.8pt,
    arc=2mm
}

\newtcolorbox{important}{
    enhanced,
    breakable,
    title={Importante},
    fonttitle=\bfseries,
    colback=yellow!8,
    colframe=orange!70!black,
    boxrule=0.8pt,
    arc=2mm
}

\newtcolorbox{example}[1]{
    enhanced,
    breakable,
    title={Ejemplo: #1},
    fonttitle=\bfseries,
    colback=green!4,
    colframe=green!40!black,
    boxrule=0.8pt,
    arc=2mm
}

\newtcolorbox{formula}[1]{
    enhanced,
    breakable,
    title={F\'ormula / Regla: #1},
    fonttitle=\bfseries,
    colback=purple!4,
    colframe=purple!50!black,
    boxrule=0.8pt,
    arc=2mm
}

\newtcolorbox{exercise}[1]{
    enhanced,
    breakable,
    title={Ejercicio: #1},
    fonttitle=\bfseries,
    colback=gray!5,
    colframe=gray!60!black,
    boxrule=0.8pt,
    arc=2mm
}

\newtcolorbox{note}{
    enhanced,
    breakable,
    title={Nota},
    fonttitle=\bfseries,
    colback=white,
    colframe=gray!50,
    boxrule=0.6pt,
    arc=2mm
}

% =====================================================
% CONFIGURACI\'ON FINAL: ENCABEZADOS, HIPERV\'INCULOS
% =====================================================
\usepackage[
    colorlinks=true,
    linkcolor=blue!50!black,
    citecolor=green!40!black,
    urlcolor=blue!60!black,
    pdftitle={\titulo},
    pdfauthor={\autor},
    pdfsubject={Apuntes universitarios},
    bookmarks=true
]{hyperref}

\pagestyle{fancy}
\fancyhf{}

\fancyhead[L]{\nouppercase{\leftmark}}
\fancyhead[R]{\materia}

\fancyfoot[C]{\thepage}

\renewcommand{\headrulewidth}{0.4pt}
\renewcommand{\footrulewidth}{0pt}

\fancypagestyle{plain}{
    \fancyhf{}
    \fancyfoot[C]{\thepage}
    \renewcommand{\headrulewidth}{0pt}
}

% =====================================================
% DOCUMENT
% =====================================================
\begin{document}

% -----------------------------------------------------
% PORTADA
% -----------------------------------------------------
\begin{titlepage}

\thispagestyle{empty}

\begin{center}

\vspace*{1cm}

{\large \textsc{\universidad}}

\vspace{0.2cm}

{\large \textsc{\facultad}}

\vspace{3cm}

{\Huge \textbf{\materia}}

\vspace{1cm}

{\LARGE \titulo}

\vspace{0.8cm}

\textit{Comprender los conceptos es el primer paso para convertir el estudio en conocimiento.}

\vspace{1.2cm}

\rule{0.70\textwidth}{0.5pt}

\vspace{0.5cm}

{\large Apuntes de estudio}

\vspace{0.5cm}

\rule{0.70\textwidth}{0.5pt}

\vfill

{\large \autor}

\vspace{0.3cm}

{\large \anio}

\end{center}

\vfill

\begin{flushleft}
\textit{Este es un modesto aporte a los estudiantes de ``Derecho Procesal Civil y Comercial'' no pretende sustituir el material oficial de la materia.}
\end{flushleft}

\end{titlepage}

% -----------------------------------------------------
% \'INDICE
% -----------------------------------------------------
\tableofcontents

% =====================================================
% CONTENIDO
% =====================================================
\chapter{Ejecuci\'on de Sentencia}

\begin{note}
Esta unidad aborda la ejecuci\'on de sentencia en el proceso civil y comercial argentino (Libro Tercero del C\'odigo Procesal Civil y Comercial de la Naci\'on, CPCCN). Se estudia su naturaleza jur\'idica como etapa del proceso de conocimiento, sus diferencias con el juicio ejecutivo, los requisitos para ejecutar una sentencia (que debe estar consentida o ejecutoriada), el procedimiento seg\'un el tipo de obligaci\'on (dar dinero l\'iquido, il\'iquido, hacer, no hacer o dar cosas), las excepciones admisibles conforme al art\'iculo 506 del CPCCN, y los t\'itulos que habilitan este proceso especial de ejecuci\'on.
\end{note}

\textit{Hilo conductor de la clase: as\'i como en un partido de f\'utbol primero se define el reglamento, luego se juega y finalmente se ejecuta el resultado (el equipo ganador cobra el trofeo), el proceso judicial tiene una etapa de conocimiento donde se discute y un fallo que solo cobra vida \'util cuando se ejecuta. Este cap\'itulo recorre ese \'ultimo tramo: c\'omo hacer efectiva la sentencia.}

\begin{table}[H]
\centering
\small
\begin{tabularx}{\textwidth}{>{\raggedright\arraybackslash}p{0.22\textwidth} X}
\toprule
\textbf{Bloque tem\'atico} & \textbf{Contenido} \\
\midrule
1. Naturaleza y encuadre & Ubicaci\'on en el C\'odigo; la ejecuci\'on como etapa del proceso de conocimiento; diferencias con el juicio ejecutivo; la sentencia como t\'itulo ejecutorio. \\
2. Requisitos para ejecutar & Sentencia consentida y ejecutoriada; plazo vencido e incumplimiento; ejecuci\'on provisoria; ejecuci\'on parcial. \\
3. Competencia y procedimiento general & Competencia del juez; condena a cantidad l\'iquida; traba de embargo; citaci\'on de venta y excepciones del art. 506; sentencia de venta. \\
4. Condena a cantidad il\'iquida & Pr\'actica de la liquidaci\'on; traslado e impugnaci\'on; criterios doctrinarios sobre la apelaci\'on. \\
5. Obligaciones de hacer, no hacer y dar cosas & Opciones del ejecutante; escritura p\'ublica; reposici\'on al estado anterior; desapoderamiento; liquidaciones especiales; otros t\'itulos ejecutables. \\
\bottomrule
\end{tabularx}
\caption{Temario general de la unidad}
\end{table}

\section{Naturaleza jur\'idica y encuadre de la ejecuci\'on de sentencia}

\subsection{Ubicaci\'on normativa en el C\'odigo Procesal}

La ejecuci\'on de sentencia se encuentra regulada dentro del Libro Tercero del C\'odigo Procesal Civil y Comercial de la Naci\'on (CPCCN), que trata los procesos de ejecuci\'on. Dentro de este libro, el T\'itulo Primero regula la ejecuci\'on de sentencias, el T\'itulo Segundo el juicio ejecutivo, y el T\'itulo Tercero las ejecuciones especiales (hipotecaria, prendaria, fiscal y comercial).

\begin{table}[H]
\centering
\begin{tabularx}{\textwidth}{>{\raggedright\arraybackslash}p{0.2\textwidth} X}
\toprule
\multicolumn{2}{l}{\textbf{Libro III del CPCCN}} \\
\midrule
T\'itulo I & Ejecuci\'on de sentencias \\
T\'itulo II & Juicio ejecutivo \\
T\'itulo III & Ejecuciones especiales (hipotecaria, prendaria, fiscal y comercial) \\
\bottomrule
\end{tabularx}
\end{table}

\subsection{La ejecuci\'on de sentencia no es un proceso independiente}

Si bien el C\'odigo incluye a la ejecuci\'on de sentencia como un proceso de ejecuci\'on m\'as, consider\'andola para alguna doctrina un proceso independiente del principal en el cual fue dictada, la posici\'on sostenida en la c\'atedra entiende que esto no es as\'i: la ejecuci\'on de sentencia es tan solo una etapa m\'as dentro del proceso de conocimiento en el cual fue dictada.

Se recuerda el esquema general del proceso de conocimiento estudiado a lo largo del curso:

\begin{itemize}
    \item \textbf{Etapa introductoria}: demanda, contestaci\'on, reconvenci\'on, excepciones.
    \item \textbf{Etapa probatoria}: producci\'on de los medios de prueba para formar convicci\'on en el juez.
    \item \textbf{Etapa decisoria}: valoraci\'on de la prueba y dictado de la sentencia.
    \item \textbf{Etapa recursiva}: mecanismos impugnativos contra la sentencia.
\end{itemize}

Una vez agotada esta \'ultima etapa, se accede a la ejecuci\'on. Aunque el propio C\'odigo trata a la ejecuci\'on de sentencia como un proceso de ejecuci\'on, para la posici\'on adoptada en la c\'atedra se trata de una etapa m\'as dentro del proceso de conocimiento.

\subsection{Diferencias entre ejecuci\'on de sentencia y juicio ejecutivo}

La intimaci\'on de pago no constituye un requisito esencial en la ejecuci\'on de sentencia, a diferencia de lo que ocurre en el juicio ejecutivo, que s\'i es un proceso de ejecuci\'on independiente. La ejecuci\'on de sentencia no requiere intimaci\'on de pago porque esa funci\'on ya queda cubierta con la notificaci\'on de la sentencia al demandado.

En cambio, el juicio ejecutivo exige la intimaci\'on de pago porque es precisamente ese acto el que habilita al ejecutado a oponer las excepciones del art\'iculo 544, a trav\'es de las cuales puede cuestionar el t\'itulo base de la ejecuci\'on.

\begin{formula}{Art. 544 CPCCN}
Excepciones admisibles en el juicio ejecutivo. El ejecutado puede oponer, al ser intimado de pago, estas excepciones para cuestionar el t\'itulo base de la ejecuci\'on: falsedad e inhabilidad de t\'itulo, entre otras.
\end{formula}

Ambas figuras tienen en com\'un que en ambas se ejecutan t\'itulos, pero con notables diferencias. En el juicio ejecutivo se ejecutan t\'itulos ejecutivos que, si bien documentan derechos, presumiblemente no tienen la certeza que posee el t\'itulo ejecutorio, que es la sentencia.

La sentencia tambi\'en constituye un t\'itulo, pero un t\'itulo que contiene un \textbf{derecho cierto}, porque ese derecho ha sido discutido con la suficiente amplitud en el juicio de conocimiento donde se dict\'o esa sentencia, y la cosa juzgada impide la discusi\'on de hechos anteriores a ella.

\begin{table}[H]
\centering
\small
\begin{tabularx}{\textwidth}{>{\raggedright\arraybackslash}p{0.2\textwidth} >{\raggedright\arraybackslash}X >{\raggedright\arraybackslash}X}
\toprule
\textbf{Aspecto} & \textbf{Juicio ejecutivo} & \textbf{Ejecuci\'on de sentencia} \\
\midrule
Tipo de t\'itulo & T\'itulo ejecutivo (derecho presunto) & T\'itulo ejecutorio: sentencia (derecho cierto) \\
Intimaci\'on de pago & Esencial. Habilita las excepciones del art. 544 & No es requisito; la sentencia ya fue notificada \\
Obligaci\'on que instrumenta & Solo obligaci\'on de dar cantidad de dinero l\'iquida o f\'acilmente liquidable & Obligaciones de dar, hacer o no hacer \\
Embargo & Optativo (puede trabarse o no) & Absolutamente esencial para la ejecuci\'on \\
Excepciones admisibles & Art. 544 (m\'as amplias, pueden abrirse a prueba) & Art. 506 (sumamente limitadas, solo hechos posteriores a la sentencia) \\
Apelaciones durante el proceso & Concedidas con tr\'amite diferido hasta apelar la sentencia de remate & Concedidas con tr\'amite diferido hasta apelar la sentencia de venta \\
\bottomrule
\end{tabularx}
\caption{Comparaci\'on entre juicio ejecutivo y ejecuci\'on de sentencia}
\end{table}

\subsection{La sentencia como t\'itulo ejecutorio: el derecho cierto}

Quien tiene a su favor una sentencia posee un t\'itulo que instrumenta un derecho cierto. Por ello se lo denomina t\'itulo ejecutorio, y porque trae aparejada una coacci\'on inmediata para agredir el patrimonio del condenado que no cumple la sentencia, aunque con algunas defensas sumamente limitadas, que se articulan a trav\'es de la oposici\'on de excepciones que consagra el art\'iculo 506 del C\'odigo Procesal.

\begin{formula}{Art. 506 CPCCN}
Excepciones admisibles en la ejecuci\'on de sentencia. Solo se consideran leg\'itimas: 1) Falsedad de la ejecutoria; 2) Prescripci\'on de la ejecutoria; 3) Pago; 4) Quita, espera o remisi\'on. En todos los casos, los hechos deben ser posteriores al dictado de la sentencia y probarse con constancias del expediente o documentos emanados del ejecutante.
\end{formula}

\section{Requisitos para ejecutar la sentencia}

\subsection{La sentencia consentida}

La principal condici\'on que debe cumplir la sentencia para ser ejecutada, conforme al art\'iculo 499, es estar consentida o ejecutoriada.

\begin{formula}{Art. 499 CPCCN}
Requisitos para ejecutar una sentencia. La sentencia debe estar consentida o ejecutoriada. Vencido el plazo de cumplimiento fijado en la sentencia --o al d\'ia siguiente de quedar firme si el juez no fij\'o plazo especial-- y a pedido de parte, el juez procede a la ejecuci\'on.
\end{formula}

Una sentencia se encuentra consentida cuando ha sido debidamente notificada a las partes y no ha sido impugnada, es decir, cuando han vencido los plazos para impugnarla a trav\'es de los recursos correspondientes sin que se lo haya hecho; o bien cuando ha sido consentida expresamente, por ejemplo mediante un escrito manifestando la conformidad con la sentencia. El consentimiento es t\'acito cuando ha vencido el plazo para impugnarla sin que ello ocurra.

\subsection{La sentencia ejecutoriada}

La sentencia est\'a ejecutoriada cuando ha sido impugnada a trav\'es de todos los mecanismos impugnativos que admit\'ia y ya no admite nuevos recursos. Tanto si ha sido consentida como si ha sido ejecutoriada, la sentencia queda firme.

\subsection{Plazo vencido, incumplimiento e instancia de parte}

Una vez firme la sentencia y vencido el plazo para su cumplimiento --plazo que puede fijar expresamente la sentencia o que, en su defecto, se entiende exigible al d\'ia siguiente de quedar firme--, a pedido de parte el juez procede a la ejecuci\'on. Se requiere inescindiblemente el pedido de la parte vencedora para proceder a ejecutar la sentencia, en virtud del principio dispositivo.

\begin{important}
Para que pueda iniciarse la ejecuci\'on de sentencia debe hallarse la sentencia incumplida por el condenado. Si vencido el plazo la sentencia se encuentra pr\'acticamente cumplida, la ejecuci\'on promovida ser\'a rechazada por haber existido cumplimiento. Se entiende que hay cumplimiento incluso cuando este se produce en virtud de la aplicaci\'on de astreintes (multa progresiva a favor del vencedor para forzar el cumplimiento del ejecutado); en tal caso no se trata de ejecuci\'on de sentencia sino de cumplimiento de la sentencia.
\end{important}

\subsection{Excepci\'on: ejecuci\'on provisoria (sentencia no firme)}

Como excepci\'on, el C\'odigo permite ejecutar la sentencia aunque no est\'e firme, mediante la llamada ejecuci\'on provisoria, sujeta a lo que en definitiva suceda a trav\'es de las v\'ias recursivas. Esto ocurre, por ejemplo, cuando la sentencia puede ser atacada mediante un recurso de apelaci\'on concedido con efecto devolutivo. El principio general es que el recurso de apelaci\'on se concede con efecto suspensivo (la apelaci\'on suspende el cumplimiento de la sentencia), salvo que la ley establezca lo contrario y disponga el efecto devolutivo.

\begin{formula}{Art. 647 CPCCN}
Efecto del recurso en materia de alimentos. El recurso de apelaci\'on contra la sentencia que condena al pago de alimentos se concede con efecto devolutivo. Esto permite al alimentado ejecutar la sentencia de primera instancia aunque est\'e apelada, form\'andose incidente de apelaci\'on con el testimonio de la sentencia de primera instancia.
\end{formula}

El alimentado beneficiado por una sentencia que condena al alimentante a pagar alimentos podr\'a ejecutar esa sentencia aunque haya sido apelada por el alimentante, dado que ese recurso se concede con efecto devolutivo. Ello protege un derecho esencial --el derecho del alimentado-- en raz\'on de la necesidad de contar con esos alimentos para su subsistencia.

\begin{formula}{Art. 498 CPCCN}
Proceso sumar\'isimo. El recurso de apelaci\'on contra la sentencia definitiva se concede con efecto devolutivo, salvo que su cumplimiento pudiere ocasionar al condenado un perjuicio irreparable, en cuyo caso el juez podr\'a concederlo con efecto suspensivo.
\end{formula}

Lo mismo sucede en el proceso sumar\'isimo, cuyo art\'iculo 498 establece que el recurso de apelaci\'on se concede con efecto devolutivo, es decir que la sentencia se cumplir\'a aunque est\'e apelada, salvo que ello pudiera causar un perjuicio irreparable al condenado, supuesto en el que se conceder\'a con efecto suspensivo.

Tambi\'en se presenta este efecto cuando se interpone recurso extraordinario contra una sentencia de c\'amara confirmatoria de la sentencia de primera instancia: si ambas sentencias son coincidentes, aunque la de segunda instancia haya sido objeto de recurso extraordinario, se permite la ejecuci\'on de la sentencia, d\'andose fianza para responder por si la Corte revocara la sentencia a trav\'es del recurso extraordinario.

\subsection{Ejecuci\'on parcial de la sentencia}

La segunda parte del art\'iculo 499 permite ejecutar parcialmente una sentencia cuando una de sus partes se hubiera consentido, no obstante haberse interpuesto alg\'un recurso respecto de los otros rubros. No se trata de una excepci\'on a la ejecutoriedad de la sentencia no consentida, sino que, habi\'endose cuestionado una parte del fallo, la parte independiente que puede cuantificarse separadamente y que qued\'o firme habilita a llevar adelante la ejecuci\'on respecto de ese rubro.

\begin{example}{Ejecuci\'on parcial en un caso de da\'nos derivados de un accidente de tr\'ansito}
Si el juez hizo lugar a los rubros da\'no emergente, lucro cesante y da\'no moral, y la parte demandada apela \'unicamente los rubros relativos al da\'no moral y al lucro cesante, la suma fijada por da\'no emergente queda firme y puede ejecutarse parcialmente por ese rubro. Se libra un testimonio dejando constancia de que ese rubro, por ese importe, ha quedado firme, para que pueda ejecutarse en primera instancia, mientras el expediente sube a c\'amara para conocer del recurso por los dem\'as rubros cuestionados.
\end{example}

\section{Competencia y procedimiento general}

\subsection{Competencia: el juez que dict\'o la sentencia}

\begin{formula}{Art. 501 CPCCN}
Competencia para la ejecuci\'on. Inc. 1: el juez que dict\'o la sentencia es competente para entender en su ejecuci\'on. Inc. 2: si el objeto de la ejecuci\'on lo impone, puede encomendarse a un juez de otra competencia territorial la realizaci\'on de determinadas diligencias, lo que no importa una excepci\'on a las reglas de competencia. Inc. 3: es competente el juez del proceso principal cuando existan conexiones directas entre causas sucesivas.
\end{formula}

El art\'iculo 501 establece que resulta competente el juez que dict\'o la sentencia. Esto se relaciona con la idea de que el proceso principal no debe entenderse aislado de la ejecuci\'on de sentencia: esta constituye una etapa m\'as de aquel, por lo que es l\'ogico que sea competente para entender en ella el mismo juez que entendi\'o en todo el proceso hasta el dictado de la sentencia.

El inciso 2 permite que, si as\'i lo impone el objeto de la ejecuci\'on, se encomiende a un juez de otra competencia territorial una diligencia determinada para ejecutar la sentencia. Esto no constituye una excepci\'on a las reglas de competencia, sino una encomienda de un juez de una jurisdicci\'on a otro, en cuanto a una delegaci\'on de la competencia para esa diligencia espec\'ifica.

\begin{example}{Delegaci\'on de diligencias en extra\'na jurisdicci\'on}
Cuando debe otorgarse la posesi\'on de un inmueble ubicado en extra\'na jurisdicci\'on, o llevarse adelante un desalojo en extra\'na jurisdicci\'on, el juez que dict\'o la sentencia (por ejemplo, con asiento en Buenos Aires) no puede trasladarse a ejecutarla en otra ciudad; por delegaci\'on, interviene el juez con competencia territorial en el lugar del inmueble.
\end{example}

El inciso 3 del art\'iculo 501 establece que es competente el juez del principal cuando medien conexiones directas entre causas sucesivas: por ejemplo, si se inici\'o como proceso principal un desalojo y, como consecuencia de ese desalojo, se inici\'o luego un proceso de da\'nos y perjuicios derivados de esa ocupaci\'on, el juez del desalojo puede entender en ambos.

\subsection{Procedimiento: condena a pagar cantidad l\'iquida}

Corresponde diferenciar si la sentencia condena al pago de una cantidad l\'iquida o, por el contrario, al pago de una cantidad il\'iquida. La sentencia condena al pago de una cantidad l\'iquida cuando esta se encuentra expresada num\'ericamente o es f\'acilmente determinable mediante una simple operaci\'on aritm\'etica.

\begin{example}{Sentencia l\'iquida aunque no exprese el monto final}
Si la sentencia condena a pagar un capital de \$500.000 con m\'as el inter\'es del 40\% anual, se entiende que la sentencia es l\'iquida aunque no exprese su monto final, dado que puede calcularse mediante un simple c\'alculo aritm\'etico que incluso puede realizar el propio juez.
\end{example}

\subsection{Traba del embargo: aspecto central de la ejecuci\'on}

\begin{important}
El embargo es absolutamente esencial porque la ejecuci\'on de sentencia pretende, ante el incumplimiento del condenado, individualizar bienes de su patrimonio para proceder a su venta y realizaci\'on efectiva, eventualmente mediante subasta. El juez tiene imperium para hacer cumplir coercitivamente sus decisiones: el obligado que no cumple voluntariamente ser\'a embargado y, una vez afectado determinado bien de su patrimonio al cumplimiento de la sentencia mediante un embargo ejecutorio, ese bien ser\'a subastado para pagar al acreedor vencedor en el proceso.
\end{important}

\subsubsection{Traba del embargo sobre bienes inmuebles}

Para pedir la traba de un embargo sobre un bien inmueble debe acompa\'narse un informe de dominio que acredite la titularidad registral en cabeza del deudor. Acreditada la titularidad, el juzgado ordena trabar embargo por la cantidad l\'iquida que consta en la sentencia, m\'as la suma que presupueste el juzgado para responder a intereses y costas.

Si el bien se encuentra en la misma jurisdicci\'on, el embargo se diligencia por oficio firmado por el secretario dirigido al Registro de la Propiedad Inmueble. Si se encuentra en extra\'na jurisdicci\'on --por ejemplo, un inmueble en otra ciudad respecto de un embargo ordenado por un juez de la Ciudad de Buenos Aires--, se traba mediante testimonio conforme a la ley 22.172.

\begin{formula}{Ley 22.172}
Comunicaciones entre jueces de distintas jurisdicciones. Establece las formas en que se comunican los jueces de distintas jurisdicciones para llevar a cabo determinadas diligencias --en este caso, la inscripci\'on del embargo ante el Registro de la Propiedad Inmueble de otra provincia--. Ese testimonio debe ser firmado por el juez y por el secretario, llevar sello de agua y otras caracter\'isticas establecidas en la ley.
\end{formula}

\subsubsection{Traba del embargo sobre bienes muebles registrables}

Si se trata de un bien mueble registrable (por ejemplo, un autom\'ovil, una embarcaci\'on o una aeronave), se embarga mediante oficio firmado por el secretario, acompa\'nado de un informe de dominio expedido por el registro correspondiente que acredite la titularidad en cabeza del demandado; el juez ordena el embargo y libra el oficio para su inscripci\'on. Si el bien se encuentra en extra\'na jurisdicci\'on, tambi\'en se procede por testimonio conforme a la ley 22.172, firmado por el juez y el secretario, con constancia de los recaudos que exige la ley.

\subsubsection{Traba del embargo sobre bienes muebles no registrables}

Si se trata de bienes muebles no registrables --es decir, m\'aquinas y herramientas que no son de uso indispensable para la profesi\'on del deudor, dado que estas \'ultimas ser\'ian inembargables--, se embarga mediante un mandamiento de embargo que firma el secretario, coordin\'andose fecha con el oficial de justicia encargado de llevarlo a cabo. El oficial de justicia se apersona en el domicilio del deudor o donde se encuentren los bienes a embargar y traba el embargo mediante un acta.

\subsection{Efectos del embargo}

El embargo concede al embargante una prioridad como primer embargante para valerse del producido del bien embargado al momento en que se efect\'ue la subasta y se convierta en dinero. Importa una indisponibilidad, una afectaci\'on del bien al proceso.

\begin{note}
El bien embargado puede ser vendido: si el deudor lo vende indicando que est\'a embargado, un tercero puede adquirirlo, pero deber\'a soportar las consecuencias de ese embargo, ya que este persigue a la cosa embargada sin perjuicio de que su titular la transfiera. Hay una indisponibilidad en el sentido de que no puede disponerse del bien embargado como si estuviera libre de gravamenes.
\end{note}

Si se desea desapoderar directamente al deudor del bien, debe solicitarse, adem\'as del embargo, el secuestro del bien, para poder trasladarlo y sustraerlo de la esfera de dominio de su due\'no.

\subsection{Citaci\'on de venta y excepciones del art\'iculo 506}

Una vez embargado un bien y afectado al proceso, se est\'a en condiciones de continuar con la ejecuci\'on. Para ello es necesario citar de venta al deudor: aunque en la etapa de ejecuci\'on de sentencia el conocimiento del juez es m\'inimo, no se excluye un cierto grado de conocimiento a trav\'es de las excepciones que puede oponer el deudor al progreso de la ejecuci\'on. De ah\'i la necesidad de citarlo de venta.

Esta citaci\'on de venta debe practicarse con posterioridad al embargo: primero se embarga y despu\'es se cita de venta, sin perjuicio de que en la pr\'actica en ocasiones ambos actos se realizan conjuntamente --por ejemplo, mediante un mandamiento en el que se traba embargo sobre determinados bienes muebles y, al mismo tiempo, se deja constancia de que se cita de venta al deudor--.

\begin{important}
La secuencia establecida por el C\'odigo es: primero embargar determinados bienes (porque sin bienes embargados no hay nada sobre lo cual hacer efectiva la sentencia) y despu\'es citar por c\'edula de venta al deudor, para que comparezca a oponer excepciones dentro de los cinco d\'ias desde la notificaci\'on.
\end{important}

\subsection{Las excepciones del art\'iculo 506: an\'alisis detallado}

Si el ejecutado no opone excepciones, el juez dicta directamente la sentencia (o resoluci\'on) de venta mandando llevar adelante la ejecuci\'on, pas\'andose a las reglas del cumplimiento para la sentencia de remate del juicio ejecutivo (art\'iculos 559 y siguientes del C\'odigo), que se estudian en la unidad n\'umero 23.

\subsubsection{Excepci\'on de falsedad de la ejecutoria}

La falsedad de la ejecutoria supone que se est\'a ejecutando una sentencia que ha sido adulterada: porque la firma no pertenece al juez, o porque se ha enmendado --sobrerraspado, adulterado-- el monto que originalmente constaba en la sentencia. No es un supuesto com\'un, pero podr\'ia darse el caso de que un testimonio solicitado salga de la esfera del expediente y sea objeto de alguna adulteraci\'on.

No puede imponerse al ejecutado el cumplimiento de una obligaci\'on que no se corresponda con la que concretamente fue dictaminada en la sentencia: debe existir una correspondencia perfecta entre lo que se ejecuta y lo decidido por la sentencia firme.

\subsubsection{Excepci\'on de prescripci\'on de la ejecutoria}

\begin{formula}{Art. 2560 del C\'odigo Civil y Comercial de la Naci\'on}
Plazo gen\'erico de prescripci\'on: cinco (5) a\'nos, excepto que est\'e previsto uno diferente en la legislaci\'on local. Es el plazo aplicable a la prescripci\'on de la ejecutoria cuando no existe plazo espec\'ifico.
\end{formula}

Corresponde aplicar el plazo gen\'erico del art\'iculo 2560 del C\'odigo Civil y Comercial de la Naci\'on, de cinco a\'nos, dado que no existe un plazo espec\'ifico para la prescripci\'on de la ejecutoria. Anteriormente ese plazo era de diez a\'nos conforme al art\'iculo 4023 del C\'odigo de V\'elez; luego se redujo a cinco a\'nos.

A partir del d\'ia en que resulta exigible el cumplimiento de la sentencia, existen cinco a\'nos para ejecutarla; si no se realizan actos que impulsen el proceso de ejecuci\'on, opera el plazo de prescripci\'on de la ejecutoria.

\begin{note}
Aunque haya transcurrido el plazo de prescripci\'on, puede deducirse la pretensi\'on de ejecuci\'on de sentencia, porque el derecho sobre la sentencia subsiste. Si la contraparte no opone la excepci\'on de prescripci\'on de la ejecutoria dentro del plazo correspondiente (cinco d\'ias), aunque estuviera vencido el plazo de prescripci\'on liberatoria, la sentencia podr\'a ejecutarse eficazmente. Existe as\'i una diferencia entre la acci\'on y el derecho: como rige el principio dispositivo, la prescripci\'on debe ser opuesta por la parte que se beneficiar\'ia con ella y no puede ser declarada de oficio por el juez.
\end{note}

\subsubsection{Excepci\'on de pago}

Si el deudor ha cancelado la obligaci\'on que surge de la sentencia, puede defenderse acreditando ese pago mediante documentos emanados del vencedor. Todas las excepciones admisibles en la ejecuci\'on de sentencia deben fundarse en hechos posteriores a la sentencia, ya que los hechos anteriores est\'an cubiertos por la eficacia de la cosa juzgada y la preclusi\'on impide volver sobre ellos. Adem\'as, hay una limitaci\'on en cuanto a los medios de prueba: solo se admiten constancias del expediente mismo o documentos emanados del vencedor, con exclusi\'on de cualquier otro medio probatorio.

\subsubsection{Excepciones de quita, espera y remisi\'on}

El inciso cuarto del art\'iculo 506 regula las excepciones de quita, de espera o de remisi\'on.

\begin{example}{Excepci\'on de quita}
Si la sentencia condena al pago de \$600.000 por da\'nos y perjuicios y existe un documento posterior a la sentencia, emanado del vencedor, en el que se realiza una quita de \$200.000, el ejecutado puede excepcionarse acompa\'nando ese documento, y la ejecuci\'on --y eventualmente el embargo-- debe reducirse a \$400.000.
\end{example}

En cuanto a la espera, corresponde cuando la parte vencedora otorga al ejecutado un plazo de cumplimiento y, no obstante ello, ejecuta la sentencia antes del vencimiento de ese plazo. En este supuesto se hace lugar a la excepci\'on y el ejecutante debe aguardar a que transcurra el plazo otorgado para poder ejecutar la sentencia.

En cuanto a la remisi\'on, consiste en el perd\'on de la deuda. T\'ecnicamente la remisi\'on implica devolver el t\'itulo donde consta la obligaci\'on al deudor, lo cual no resulta posible en este \'ambito porque se trata de un expediente ya iniciado; si se acompa\'na el instrumento correspondiente que acredita el perd\'on de la deuda y el juez hace lugar a la excepci\'on, se ordena levantar el embargo.

\subsection{Excepciones no enumeradas en el art\'iculo 506 admitidas por la jurisprudencia}

El art\'iculo 506 dispone que solo se consideran leg\'itimas las excepciones que enumera, pero la jurisprudencia ha entendido que no pueden dejarse de lado otras excepciones, como la inhabilidad de t\'itulo: si el ejecutante pretende ejecutar a quien no es el demandado condenado, el pretendido ejecutado puede excepcionarse mediante la falta de legitimaci\'on pasiva, incluida dentro de la excepci\'on de inhabilidad de t\'itulo.

Lo mismo se ha admitido respecto de otras excepciones vinculadas a presupuestos procesales, como la excepci\'on de incompetencia: si se ejecuta la sentencia ante un juez incompetente, aunque no est\'e expresamente se\'nalada dentro de las enumeradas en el art\'iculo 506, el deudor puede oponer la excepci\'on de incompetencia.

\subsection{Tr\'amite de las excepciones y sentencia de venta}

Opuesta alguna de las excepciones por el ejecutado, corresponde dar traslado al ejecutante por cinco d\'ias para que las conteste. Contestado el traslado, el juez debe resolver.

Si el juez hace lugar a las excepciones, ordena levantar el embargo (en el caso de pago o de prescripci\'on de la ejecutoria); si es la excepci\'on de quita, adecua el embargo a la nueva suma; si es la de espera, hace lugar a la excepci\'on pero deja trabado el embargo, para que el ejecutante pueda continuar una vez verificado el plazo correspondiente.

Si en cambio el juez rechaza la excepci\'on opuesta, dicta la sentencia de venta mandando llevar adelante la ejecuci\'on. Esta resoluci\'on puede ser apelada por el ejecutado; el recurso de apelaci\'on contra la sentencia de venta se concede con efecto suspensivo y tr\'amite inmediato, pero si el ejecutante otorga fianza para responder por da\'nos y perjuicios, puede transformar el efecto de la apelaci\'on en devolutivo, ejecutando la sentencia no obstante la apelaci\'on de la parte ejecutada.

Si se ordena levantar el embargo, el ejecutante puede apelar esa resoluci\'on con efecto suspensivo y tr\'amite inmediato. Una vez firme, se contin\'ua con el tr\'amite del cumplimiento de la sentencia de remate del juicio ejecutivo, etapa com\'un a ambos procesos.

\section{Condena a cantidad il\'iquida y pr\'actica de liquidaci\'on}

\subsection{Cu\'ando corresponde practicar liquidaci\'on}

Distinto es el procedimiento cuando la sentencia no establece un monto determinado que deba abonar el condenado --por ejemplo, en concepto de intereses (se fija la tasa pero no el monto), de frutos a restituir, o de da\'nos y perjuicios cuyo derecho se reconoce al actor pero que, por distintas cuestiones procesales, el juez no puede fijar en el acto del fallo.

\begin{formula}{Art. 165 CPCCN}
Condena al pago de frutos, intereses o da\'nos. El juez, al pronunciar sentencia, debe fijar el importe de los frutos, intereses, da\'nos y perjuicios con arreglo a lo dispuesto por el art\'iculo 165. Cuando ello no sea posible, el juez debe fijar las bases para que se practique la liquidaci\'on conforme a las pautas que establezca.
\end{formula}

El art\'iculo 165 exige, en primer lugar, que el juez determine la cantidad en forma l\'iquida respecto de estos rubros, pero reconociendo que puede haber impedimentos que obstaculicen esa posibilidad; en tal caso, exige al juez que al menos fije las pautas para que la liquidaci\'on pueda establecerse con posterioridad.

\begin{important}
La liquidaci\'on debe estar en consonancia y ser correlativa con lo establecido en la sentencia: no puede incluir, al practicarse, rubros que vayan m\'as all\'a de lo solicitado en las pretensiones articuladas al iniciarse el proceso, porque de lo contrario ser\'ia nugatoria la ejecuci\'on basada en una liquidaci\'on que incluyera conceptos indebidos.
\end{important}

\subsection{Qui\'en practica la liquidaci\'on y en qu\'e plazo}

\begin{formula}{Art. 503 CPCCN (liquidaci\'on)}
Liquidaci\'on. El acreedor podr\'a practicar la liquidaci\'on dentro del plazo de diez d\'ias contados desde que la sentencia sea exigible. Dar\'a traslado a la otra parte por cinco d\'ias. Si hay conformidad, o si vencido el plazo el deudor no la impugna, el juez aprobar\'a la liquidaci\'on presentada, disponiendo que se la tenga como base para el embargo y la subasta.
\end{formula}

El C\'odigo establece que practicar la liquidaci\'on es una facultad del vencedor, para lo cual dispone de diez d\'ias contados a partir de que sea exigible el cumplimiento --es decir, desde que la sentencia qued\'o firme--. De esa liquidaci\'on debe darse traslado a la otra parte por c\'edula, por cinco d\'ias. Si la contraparte la consiente expresamente, o guarda silencio consinti\'endola t\'acitamente por no impugnarla, el juez aprueba la liquidaci\'on y se contin\'ua como en el caso de cantidad l\'iquida, trab\'andose embargo sobre el monto de la liquidaci\'on y lo que se presupueste para intereses y costas.

\subsection{Impugnaci\'on de la liquidaci\'on}

\begin{formula}{Art. 504 CPCCN}
Impugnaci\'on de la liquidaci\'on. Si hubiere impugnaci\'on a la liquidaci\'on, se formar\'a incidente siguiendo las normas de los art\'iculos 178 y siguientes. Las partes deber\'an ofrecer la prueba de la que intenten valerse en sus respectivos escritos de impugnaci\'on y contestaci\'on.
\end{formula}

Si el vencido impugna la liquidaci\'on --porque entiende que no corresponden los rubros, o porque, correspondiendo, est\'an mal o defectuosamente calculados-- se forma un incidente conforme al art\'iculo 504, siguiendo las pautas del art\'iculo 178 y siguientes. En ese incidente las partes deben ofrecer la prueba de la que intenten valerse para fundar sus posiciones. Vencidos los plazos y contestada la impugnaci\'on, el juez debe resolver aprobando la liquidaci\'on que corresponda.

\subsection{Criterios de apelaci\'on sobre la liquidaci\'on aprobada}

\begin{important}
Existen dos posiciones doctrinarias sobre el efecto con que se concede la apelaci\'on contra la resoluci\'on que aprueba la liquidaci\'on. Para Alsina, por tratarse de un incidente, el recurso de apelaci\'on debe concederse en relaci\'on, con efecto suspensivo y tr\'amite inmediato. Para Palacio, en cambio, se aplican las normas del art\'iculo 509, que dispone que todas las apelaciones en la ejecuci\'on de sentencia se conceden con tr\'amite diferido.

Esta diferencia altera sustancialmente la soluci\'on del caso: si se sigue el tr\'amite diferido, la liquidaci\'on aprobada --aunque apelada por el ejecutado-- sirve de base para trabar el embargo y continuar con la ejecuci\'on, y reci\'en podr\'a fundarse esa apelaci\'on junto con una eventual apelaci\'on contra la sentencia de venta. Si en cambio se sostiene el criterio del tr\'amite inmediato, corresponder\'ia suspender el cumplimiento de la resoluci\'on --por tener efecto suspensivo-- hasta que la c\'amara la resuelva, y reci\'en confirmada por esta, en primera instancia se proceder\'ia a trabar el embargo y continuar la ejecuci\'on.
\end{important}

\section{Obligaciones de hacer, no hacer y dar cosas}

\subsection{Condena a hacer alguna cosa}

\begin{formula}{Art. 503 CPCCN (condena a hacer)}
Si la condena es a hacer alguna cosa y el obligado no la realiza dentro del plazo correspondiente, el ejecutante puede optar por: a) que la cosa sea realizada por un tercero a costa del ejecutado; o b) que se le paguen los da\'nos y perjuicios provenientes de la inejecuci\'on.
\end{formula}

Si la condena es a hacer alguna cosa y el deudor cumple realizando lo que le fue ordenado dentro del plazo correspondiente, no hay ejecuci\'on. Si no cumple, el ejecutante puede promover la ejecuci\'on de sentencia, optando --dado que el deudor no puede ser obligado mediante la fuerza f\'isica a hacer algo contra su voluntad-- entre solicitar que la cosa sea realizada por un tercero, o que se lo obligue a resarcir los da\'nos y perjuicios provenientes de la inejecuci\'on.

Si se opta por que un tercero realice la prestaci\'on, corresponde estimar --no liquidar-- el costo de esa actividad, y dar traslado al vencido para que se expida sobre si ese valor se corresponde con el objeto de la obligaci\'on que ahora realizar\'a un tercero a su costa. La opci\'on corresponde al acreedor, pero a costa del ejecutado.

\subsection{Condena a otorgar escritura p\'ublica}

\begin{formula}{Art. 502 CPCCN}
Condena a otorgar escritura p\'ublica. Si la sentencia condena al otorgamiento de escritura p\'ublica y este no es cumplido por el ejecutado, el juez la suscribir\'a a su costa, sin perjuicio de que, si deviniera de imposible cumplimiento por alguna circunstancia, esa escrituraci\'on se transforme en la obligaci\'on de resarcir los da\'nos y perjuicios provenientes de esa imposibilidad.
\end{formula}

Si la sentencia condena al otorgamiento de escritura p\'ublica y esto no es cumplido por el ejecutado, se designa a quien la suscriba a su costa. Si deviniera de imposible cumplimiento por alguna circunstancia, la condena se transforma en la obligaci\'on de resarcir los da\'nos y perjuicios provenientes de esa imposibilidad, fij\'andose estos por el procedimiento del art\'iculo 513.

\subsection{Condena a no hacer}

\begin{formula}{Art. 513 CPCCN}
Da\'nos y perjuicios por inejecuci\'on. Los da\'nos y perjuicios derivados del incumplimiento se determinar\'an conforme a las pautas de los art\'iculos 503 y 504 (liquidaci\'on con v\'ia incidental) o bien a trav\'es de juicio ordinario o sumar\'isimo, seg\'un lo determine el juez de manera irrecurrible.
\end{formula}

Cuando se trata de una condena a no hacer y esta es incumplida por el ejecutado --es decir, hace algo contrario a lo que le hab\'ia sido ordenado--, el acreedor puede optar entre solicitar que las cosas se restituyan al estado anterior, o que ello se traduzca en la indemnizaci\'on de da\'nos y perjuicios conforme al art\'iculo 513. Estos da\'nos y perjuicios se determinan mediante el procedimiento de la liquidaci\'on con derivaci\'on a la v\'ia incidental, o bien por juicio ordinario o sumar\'isimo, seg\'un lo determine el juez de manera irrecurrible.

\subsection{Condena a dar cosas: mandamiento de desapoderamiento}

Si la sentencia no es cumplida voluntariamente por el ejecutado, pueden aplicarse astreintes --multas progresivas conminatorias a favor de la otra parte--. Si la obligaci\'on de dar alguna cosa no se cumple voluntariamente, el vencedor puede ejecutar la sentencia mediante un mandamiento que ordena librar el juez para desapoderar --aun contra la voluntad del deudor-- al vencido de aquello que fue ordenado restituir por la sentencia.

Con ese mandamiento se cita, en el mismo acto en que se diligencia, al ejecutado para que oponga las excepciones del art\'iculo 506.

\begin{note}
Si por alguna circunstancia no puede cumplirse la entrega de la cosa --porque ha desaparecido o su cumplimiento resulta imposible--, el vencedor puede solicitar que se compense con una suma equivalente al valor de la cosa. La determinaci\'on de dicho valor se realiza tambi\'en conforme a las pautas de los art\'iculos 503 y 504, mediante la v\'ia incidental o por juicio ordinario o sumar\'isimo, seg\'un lo establezca el juez, de manera irrecurrible.
\end{note}

\subsection{Liquidaciones especiales}

\begin{formula}{Art. 516 CPCCN}
Liquidaciones especiales. Cuando las liquidaciones impliquen cuentas complicadas de lenta o dif\'icil justificaci\'on o que requieran conocimientos especiales, podr\'an someterse a la decisi\'on de peritos \'arbitros o amigables componedores. El laudo que emitan ser\'a vinculante y obligatorio para las partes. Lo mismo se aplica para liquidaciones de sociedades conyugales o determinaci\'on del car\'acter propio o ganancial de bienes.
\end{formula}

Cuando las liquidaciones impliquen cuentas complicadas de lenta o dif\'icil justificaci\'on, o que requieran conocimientos especiales, puede dejarse su determinaci\'on librada a la decisi\'on de peritos \'arbitros o amigables componedores. Tanto el perito \'arbitro como el amigable componedor emiten un laudo vinculante y obligatorio para las partes, que el juez debe respetar en esas cuestiones t\'ecnicas.

\subsection{Otros t\'itulos ejecutables por este procedimiento}

\begin{formula}{Art. 500 CPCCN}
T\'itulos ejecutables por el procedimiento de ejecuci\'on de sentencia: 1) Las sentencias de los tribunales judiciales y los laudos de tribunales arbitrales (aunque estos carecen de imperium); 2) Las transacciones o acuerdos homologados; 3) (Inciso incorporado por la ley 26.589) Los acuerdos celebrados en mediaci\'on con la firma certificada del mediador, salvo que incluyan menores o incapaces, en cuyo caso requieren homologaci\'on judicial; 4) Las multas procesales; 5) Los honorarios regulados judicialmente.
\end{formula}

El art\'iculo 500 establece que, adem\'as de aplicarse el procedimiento de ejecuci\'on de sentencia respecto de las dictadas por tribunales judiciales o de los laudos de tribunales arbitrales --que, si bien tienen jurisdicci\'on, carecen de imperium para hacer cumplir coercitivamente sus decisiones--, ese procedimiento tambi\'en se aplica a las transacciones o acuerdos homologados.

Las transacciones son aquellas mediante las cuales las partes extinguen obligaciones litigiosas o dudosas, y los acuerdos homologados pueden provenir, por ejemplo, de una conciliaci\'on llevada a cabo por el juez en el \'ambito de la audiencia preliminar.

\begin{note}
Los acuerdos conciliatorios celebrados ante el mismo juez en la audiencia preliminar requieren homologaci\'on. En cambio, el inciso cuarto del art\'iculo 500, incorporado por la ley 26.589, permite ejecutar por este procedimiento un acuerdo celebrado en mediaci\'on con la firma certificada del mediador, sin necesidad de homologaci\'on --salvo que se trate de acuerdos que incluyan menores o incapaces, caso en el cual su representante debe solicitar al juez la homologaci\'on con intervenci\'on del Ministerio Pupilar--.
\end{note}

La importancia del acuerdo celebrado en mediaci\'on radica en que su incumplimiento permite sortear todas las etapas del proceso de conocimiento y acceder directamente a la \'ultima etapa --la ejecuci\'on de sentencia-- para exigir su cumplimiento conforme a las pautas estudiadas.

El art\'iculo 500, en su inciso segundo, tambi\'en prev\'e la ejecuci\'on de las multas procesales --por ejemplo, las impuestas por temeridad y malicia a favor de una de las partes--, y finalmente corresponde que se ejecute por este procedimiento el cobro de honorarios regulados judicialmente.

\subsection*{Reflexi\'on final}

El C\'odigo establece un procedimiento de ejecuci\'on de sentencia abreviado, limitado en cuanto a la apelaci\'on y en cuanto a las defensas admisibles, dada la trascendencia del t\'itulo ejecutorio que, como se se\'nal\'o al principio, instrumenta un derecho cierto. Ello permite transitar una etapa con un m\'inimo de conocimiento y una m\'inima posibilidad de prueba de parte del ejecutado, y hace posible efectivizar la sentencia y concretarla pr\'acticamente, de modo de satisfacer el derecho que pretend\'ia el actor al iniciar la demanda: la sentencia no es una mera declaraci\'on de certeza respecto de ese derecho, sino que la pretensi\'on requiere su efectivo cumplimiento.

% TODO: la clase cierra con una reflexi\'on atribuida gen\'ericamente a "alg\'un autor" sin identificarlo en las notas originales: para tener derecho hace falta contar con la norma que lo acuerde, saberlo exponer, que el juez quiera reconocerlo y que el deudor pueda pagarlo.

% =====================================================
% CORNELL
% =====================================================
\chapter*{Cuadro Cornell de repaso}
\addcontentsline{toc}{chapter}{Cuadro Cornell de repaso}

\begin{longtable}{
    >{\raggedright\arraybackslash}p{0.30\textwidth}
    >{\raggedright\arraybackslash}p{0.60\textwidth}
}
\toprule
\textbf{Preguntas / palabras clave} & \textbf{Notas / respuestas} \\
\midrule
\endfirsthead
\toprule
\textbf{Preguntas / palabras clave} & \textbf{Notas / respuestas} \\
\midrule
\endhead
\bottomrule
\endfoot
\bottomrule
\endlastfoot

\textbf{¿Por qu\'e la ejecuci\'on de sentencia NO es un proceso independiente?} &
Porque constituye la culminaci\'on l\'ogica del proceso de conocimiento: est\'a \'intimamente ligada a la sentencia que se ejecuta, el mismo juez que conoci\'o en la causa es competente, y la intimaci\'on de pago no es un requisito esencial (ya se notifica en la sentencia). El C\'odigo la trata como proceso independiente por razones sistem\'aticas, pero la posici\'on de la c\'atedra y el propio articulado del C\'odigo la ubican como una etapa m\'as. \\

\textbf{¿Qu\'e diferencia al t\'itulo ejecutorio del t\'itulo ejecutivo?} &
El t\'itulo ejecutivo (juicio ejecutivo) contiene un derecho presunto, discutible con mayor amplitud (art. 544). El t\'itulo ejecutorio (sentencia) contiene un derecho cierto porque ya fue debatido exhaustivamente en el proceso de conocimiento. La cosa juzgada impide discutir hechos anteriores al dictado de la sentencia. \\

\textbf{¿Cu\'ando se considera consentida una sentencia?} &
Cuando: a) fue debidamente notificada y no fue impugnada en el plazo; b) fue consentida expresamente por escrito; c) fue apelada pero el recurso se abandon\'o. El consentimiento puede ser expreso o t\'acito. \\

\textbf{¿Qu\'e significa que la sentencia est\'e ejecutoriada?} &
Que fue impugnada por todos los recursos que admit\'ia y ya no admite nuevos embates: se agotaron todos los mecanismos impugnativos disponibles. \\

\textbf{¿Cu\'ales son los requisitos para ejecutar una sentencia?} &
1) Sentencia consentida o ejecutoriada. 2) Vencido el plazo de cumplimiento (fijado por el juez o al d\'ia siguiente de quedar firme). 3) Incumplimiento del condenado. 4) Instancia de parte (petici\'on del vencedor). Rige el principio dispositivo: sin pedido de parte no hay ejecuci\'on. \\

\textbf{¿Cu\'ando se puede ejecutar una sentencia que NO est\'a firme?} &
Excepcionalmente, cuando la apelaci\'on se concede con efecto devolutivo: a) Alimentos (art. 647 CPCCN); b) Proceso sumar\'isimo (art. 498, salvo perjuicio irreparable); c) Recurso extraordinario ante sentencias de primera y segunda instancia coincidentes. En todos los casos puede exigirse fianza para transformar el efecto suspensivo en devolutivo. \\

\textbf{¿Cu\'al es la diferencia entre ejecuci\'on total y ejecuci\'on parcial?} &
La ejecuci\'on parcial procede cuando solo una parte del fallo qued\'o firme (el resto fue apelado). Se libra testimonio acreditando que ese rubro espec\'ifico qued\'o consentido, se ejecuta en primera instancia ese rubro, y el expediente sube a c\'amara por los rubros cuestionados. \\

\textbf{¿Por qu\'e el embargo es esencial en la ejecuci\'on de sentencia?} &
Porque el procedimiento tiene por finalidad afectar bienes del deudor para su posterior subasta: sin bien afectado al proceso, no puede continuarse la ejecuci\'on. En el juicio ejecutivo el embargo es optativo; en la ejecuci\'on de sentencia es absolutamente esencial. \\

\textbf{¿C\'omo se traba el embargo seg\'un el tipo de bien?} &
Inmueble (misma jurisdicci\'on): oficio del secretario al Registro de la Propiedad. Inmueble (otra jurisdicci\'on): testimonio ley 22.172, firmado por juez y secretario con sello de agua. Mueble registrable: oficio del secretario m\'as informe de dominio del registro correspondiente. Mueble no registrable: mandamiento de embargo firmado por el secretario, diligenciado por el oficial de justicia. \\

\textbf{¿Qu\'e es la citaci\'on de venta y cu\'ando se produce?} &
Es el acto procesal que habilita al ejecutado a oponer excepciones. Se practica despu\'es del embargo (aunque en la pr\'actica puede hacerse en el mismo mandamiento). El ejecutado tiene cinco d\'ias para oponer las excepciones del art. 506 CPCCN. \\

\textbf{¿Cu\'ales son las excepciones del art. 506 CPCCN?} &
1) Falsedad de la ejecutoria (adulteraci\'on de la sentencia). 2) Prescripci\'on de la ejecutoria (plazo de 5 a\'nos, art. 2560 CCCN). 3) Pago (documentado con recibo emanado del vencedor). 4) Quita, espera o remisi\'on (documentos posteriores a la sentencia, emanados del ejecutante). Condici\'on com\'un: hechos posteriores a la sentencia; prueba limitada a constancias del expediente o documentos del ejecutante. \\

\textbf{¿Qu\'e pasa si el ejecutado NO opone excepciones?} &
El juez dicta directamente la sentencia de venta mandando llevar adelante la ejecuci\'on. A partir de all\'i se aplican las reglas del cumplimiento para la sentencia de remate del juicio ejecutivo (art. 559 y siguientes CPCCN, tratados en la unidad 23). \\

\textbf{¿Qu\'e excepciones admite la jurisprudencia m\'as all\'a del art. 506?} &
La jurisprudencia ampl\'ia las excepciones para incluir: a) inhabilidad de t\'itulo (falta de legitimaci\'on pasiva); b) incompetencia (ejecutar ante juez incompetente). La enumeraci\'on del art. 506 no es estrictamente taxativa seg\'un la jurisprudencia predominante. \\

\textbf{¿Qu\'e sucede cuando la sentencia condena a una cantidad il\'iquida?} &
Corresponde practicar liquidaci\'on (art. 503 CPCCN). El vencedor tiene diez d\'ias desde que la sentencia es exigible. Se da traslado a la otra parte por cinco d\'ias. Si hay conformidad o silencio, el juez aprueba la liquidaci\'on y se sigue con el tr\'amite de embargo. Si hay impugnaci\'on, se forma incidente (art. 504, reglas del art. 178 y siguientes). \\

\textbf{¿Qu\'e ocurre con la condena a HACER algo?} &
El ejecutante puede optar (art. 503): a) que un tercero realice la cosa a costa del ejecutado (se estima el costo y se da traslado); b) pago de da\'nos y perjuicios por inejecuci\'on (se cuantifican por liquidaci\'on/incidente o juicio ordinario/sumar\'isimo seg\'un determine el juez, irrecurriblemente). \\

\textbf{¿Qu\'e pasa si la condena es a NO HACER algo?} &
El acreedor puede optar (art. 513): a) restituci\'on al estado anterior; b) indemnizaci\'on de da\'nos y perjuicios, cuantificados por liquidaci\'on/incidente o juicio ordinario/sumar\'isimo, seg\'un determine el juez de manera irrecurrible. \\

\textbf{¿Qu\'e t\'itulos habilitan la ejecuci\'on de sentencia adem\'as de la sentencia judicial?} &
Art. 500 CPCCN: 1) laudos arbitrales (el \'arbitro carece de imperium para ejecutar); 2) transacciones y acuerdos homologados; 3) acuerdos de mediaci\'on con firma certificada del mediador (ley 26.589), salvo menores o incapaces, que requieren homologaci\'on; 4) multas procesales; 5) honorarios regulados judicialmente. \\

\textbf{¿Cu\'al es la paradoja del art. 500 entre acuerdo judicial y acuerdo de mediaci\'on?} &
El acuerdo conciliatorio celebrado ante el mismo juez en la audiencia preliminar requiere homologaci\'on judicial para ser ejecutable, mientras que el acuerdo de mediaci\'on con firma certificada del mediador no la requiere (salvo menores o incapaces). La ley 26.589 introdujo esta distinci\'on al incorporar el inciso 4 del art. 500. \\

\midrule
\multicolumn{2}{p{0.92\textwidth}}{
\textbf{Resumen:} La ejecuci\'on de sentencia es la etapa final del proceso de conocimiento --no un proceso aut\'onomo--, mediante la cual el Estado pone su imperium al servicio del vencedor para hacer efectivo un derecho declarado cierto por sentencia firme. A diferencia del juicio ejecutivo, que ejecuta t\'itulos con derechos presuntos y admite defensas m\'as amplias (art. 544 CPCCN), la ejecuci\'on de sentencia opera sobre un t\'itulo ejecutorio que ya super\'o el debate exhaustivo del proceso de conocimiento: la cosa juzgada material impide discutir hechos anteriores. Por eso las excepciones admisibles (art. 506 CPCCN) se reducen a falsedad de la ejecutoria, prescripci\'on, pago y quita/espera/remisi\'on --todas fundadas en hechos posteriores a la sentencia y probables solo con documentos del expediente o del ejecutante--. El procedimiento requiere que la sentencia est\'e consentida o ejecutoriada, que haya vencido el plazo de cumplimiento y que exista incumplimiento del condenado, todo a instancia de parte. El embargo es esencial: traba bienes para su posterior subasta. Seg\'un el tipo de obligaci\'on condenada (dar dinero, hacer, no hacer, dar cosas), el proceso adopta distintos carriles: liquidaci\'on (si la condena es il\'iquida), tercer\'ismo a costa del deudor o indemnizaci\'on (si la condena es a hacer o no hacer), y mandamiento de desapoderamiento (si es a dar cosas). Otros t\'itulos equiparados a la sentencia --laudos arbitrales, transacciones homologadas, acuerdos de mediaci\'on certificados, multas procesales y honorarios-- tambi\'en se ejecutan por este procedimiento, lo que subraya su car\'acter de herramienta central en la tutela jurisdiccional efectiva.
} \\
\bottomrule

\end{longtable}

\end{document}
